\begin{tikzpicture}[
    >=Latex,
    every node/.style={font=\small},
]
    % =========================================================
    % (a) Persistence-aehnliches Diagramm
    % =========================================================
    \node[font=\bfseries, anchor=west] at (0, 3.0)
        {(a) Topologie-Evolution über die Iso-Stufen};

    % Achsen
    \draw[->, thick] (-0.3, 0) -- (10, 0)
        node[anchor=west, font=\scriptsize] {$\varphi^*$};

    % X-Skala
    \foreach \x/\v in {0/0, 2.5/0{,}25, 5/0{,}5, 7.5/0{,}75, 10/1} {
        \draw (\x, -0.1) -- (\x, 0.1);
        \node[anchor=north, font=\scriptsize] at (\x, -0.15) {\v};
    }

    % Iso-Stufen (gestrichelte Hilfslinien)
    \foreach \x in {1, 2, 3, 4, 5, 6, 7, 8, 9} {
        \draw[gray!40, dotted, very thin] (\x, -1.8) -- (\x, 2.5);
    }

    % --- LINIEN (zuerst zeichnen, damit Symbole obenauf liegen) ---
    % Stamm-Branch (b=0) von phi*=0 bis exakt phi*=0.5 (x=5)
    \draw[ultra thick, blue!70!black] (0.3, 0) -- (5.0, 0);
    
    % Zweig 1 (b=1, oben) nach Split
    \draw[ultra thick, orange!85!black] (5.0, 0) -- (5.5, 1.5) -- (10, 1.5);
    
    % Zweig 2 (b=2, unten) nach Split
    \draw[ultra thick, green!50!black] (5.0, 0) -- (5.5, -1.5) -- (10, -1.5);


    % --- SYMBOLE & TEXTE (liegen nun sauber über den Linien) ---
    % Stamm Text
    \node[font=\scriptsize, blue!70!black, anchor=south, align=center]
        at (2.5, 0.15) {Stamm $b=0$};

    % Zweig Texte
    \node[font=\scriptsize, orange!85!black, anchor=south]
        at (8, 1.65) {Zweig $b=1$};
    \node[font=\scriptsize, green!50!black, anchor=north]
        at (8, -1.65) {Zweig $b=2$};

    % Birth-Symbol (offener Kreis) am Stammbeginn
    \draw[blue!70!black, fill=white, thick] (0.3, 0) circle (0.13);
    \node[font=\scriptsize, anchor=south] at (0.3, 0.25) {Birth};

    \draw[red!80!black, fill=red!20, thick]
        (5.0, -0.18) -- (5.0, 0.18) -- (5.45, 0) -- cycle;
    \node[font=\scriptsize, red!80!black, anchor=south]
        at (4.5, 0.15) {Split};

    % Death-Symbole (gefuellte Kreise) am Ende
    \draw[orange!85!black, fill=orange!85!black, thick]
        (10, 1.5) circle (0.13);
    \draw[green!50!black, fill=green!50!black, thick]
        (10, -1.5) circle (0.13);
    \node[font=\scriptsize, anchor=west] at (10.2, 1.5) {Death};
    \node[font=\scriptsize, anchor=west] at (10.2, -1.5) {Death};

    % =========================================================
    % (b) T-Traeger-Referenz (klein, unten links)
    % =========================================================
    \begin{scope}[xshift=1.5cm, yshift=-4.2cm, scale=0.6]
        \node[font=\bfseries, anchor=west] at (-1.5, 2.2)
            {(b) Zugehöriges Bauteil};

        % Farbige Füllungen
        \fill[blue!30] (-0.3, -0.5) rectangle (0.3, 1.5);     % Kompletter Stamm
        \fill[green!40] (-1.5, 1.0) rectangle (-0.3, 1.5);    % Linker Querbalken (b=2)
        \fill[orange!40] (0.3, 1.0) rectangle (1.5, 1.5);     % Rechter Querbalken (b=1)

        \draw[thick] (-0.3, 1.0) -- (-0.3, 1.5);
        \draw[thick] (0.3, 1.0) -- (0.3, 1.5);

        % Äußere Umrandung
        \draw[thick]
            (-0.3, -0.5) -- (0.3, -0.5) -- (0.3, 1.0) --
            (1.5, 1.0) -- (1.5, 1.5) --
            (-1.5, 1.5) -- (-1.5, 1.0) --
            (-0.3, 1.0) -- cycle;

        % Branch-Beschriftungen
        \node[font=\tiny, blue!70!black] at (0, -1.0) {$b=0$};
        \node[font=\tiny, orange!85!black] at (0.9, 0.75) {$b=1$};
        \node[font=\tiny, green!50!black] at (-0.9, 0.75) {$b=2$};
    \end{scope}

    % =========================================================
    % (c) Legende der Topologie-Ereignisse (rechts neben Bauteil)
    % =========================================================
    \begin{scope}[xshift=6.5cm, yshift=-4.2cm]
        \node[font=\bfseries, anchor=west] at (0, 1.5)
            {(c) Topologie-Ereignisse};

        % Birth
        \draw[black, fill=white, thick] (0.2, 1.0) circle (0.10);
        \node[font=\scriptsize, anchor=west] at (0.4, 1.0)
            {Birth: neue Komponente};

        % Death
        \draw[black, fill=black, thick] (0.2, 0.6) circle (0.10);
        \node[font=\scriptsize, anchor=west] at (0.4, 0.6)
            {Death: Komponente endet};

        % Split
        \draw[red!80!black, fill=red!20, thick]
            (0.1, 0.1) -- (0.1, 0.3) -- (0.3, 0.2) -- cycle;
        \node[font=\scriptsize, anchor=west] at (0.4, 0.2)
            {Split: Aufspaltung};

        % Merge
        \draw[red!80!black, fill=red!20, thick]
            (0.3, -0.4) -- (0.3, -0.2) -- (0.1, -0.3) -- cycle;
        \node[font=\scriptsize, anchor=west] at (0.4, -0.3)
            {Merge: Zusammenführung};
    \end{scope}

\end{tikzpicture}